% Volume VI: Machine Learning Mathematics
% Continuity note for Chapter 31.
% Previous dependencies: Chapters 1-2 provide typed functions and domains; Chapters 5-7 provide least squares, projections, matrix calculus, and shape discipline; Chapters 13-18 provide probability, expectation, likelihood, Bayes, risk, and information-theoretic losses.
% New durable concepts: supervised learning as statistical function approximation; input, output, action, hypothesis class, loss, learner, training sample, population risk, empirical risk, empirical risk minimization, likelihood-loss links, train/test risk, generalization gap, bias-variance decomposition, finite-class generalization bounds, no-free-lunch intuition, and inductive bias.
% Forward hooks: Chapter 32 specializes the learning setup to linear models, GLMs, and classifiers; Chapter 33 adds regularization, kernels, RKHS methods, and Gaussian-process priors; Chapter 34 replaces direct prediction functions by latent-variable generative models; later deep-learning and neuroscience chapters reuse risk, empirical risk, validation, and inductive bias.
% Notation commitments: Inputs are x in \mathcal X, outputs are y in \mathcal Y, predictions/actions are a in \mathcal A, data pairs are (X,Y) distributed according to P, samples are S_n=((X_i,Y_i))_{i=1}^n, functions are f:\mathcal X\to\mathcal A, hypothesis classes are \mathcal F, losses are \ell:\mathcal A\times\mathcal Y\to[0,\infty), population risk is R(f), empirical risk is \widehat R_n(f), and learned functions are \widehat f.
% Figure plan: no missing images are included. Proposed ImageGen prompts are recorded in imagegen-figures/ch31-figure-metadata.md.
\section{Chapter 31: Learning as Statistical Function Approximation}
\label{ch:learning-statistical-function-approximation}

The previous volumes built the mathematical language needed to describe data, uncertainty, approximation, and optimization. Machine learning begins when those ideas are combined into a single statistical question: from finitely many examples, choose a function that will make good predictions on examples not yet seen.

The central problem of this chapter is not merely how to fit data. It is how to define the prediction task, choose a function class, measure error, replace an unknown population objective by a computable sample objective, and understand why fitting the sample may or may not generalize. This is the statistical foundation underneath regression, classification, neural decoding, tuning-curve estimation, and modern deep learning.

\paragraph{Chapter problem.}
How can learning be formulated as choosing a function from finite data so that its expected loss on future data is small?

\paragraph{Section map.}
Section 31.1 defines the formal learning problem: data distribution, hypothesis class, loss, risk, and learner. Section 31.2 studies empirical risk minimization and the link between losses and likelihoods. Section 31.3 explains generalization through train/test risk, bias-variance intuition, and a finite-class bound. Section 31.4 states the no-free-lunch lesson and shows why inductive bias is not optional.

\subsection{31.1 Formal learning problem}

\paragraph{Guiding question.}
What are the mathematical objects in a supervised learning problem, and what function is learning trying to approximate?

\paragraph{Beginner-facing intuition.}
A supervised learning problem starts with examples. Each example has an input and a desired output: an image and a class label, a vector of spike counts and a reach direction, a stimulus and a neural response, a sentence prefix and the next token. The learner does not see the whole population of possible examples. It sees a finite sample and must choose a prediction rule.

The important shift is from memorizing the table to approximating a function. If similar inputs tend to have similar outputs, or if the output depends on a small set of features, a function fitted on training examples can predict future examples. If there is no stable relation between inputs and outputs, learning has nothing to exploit.

\paragraph{Formal definitions.}
Let \(\mathcal X\) be an input space, \(\mathcal Y\) an output space, and \(\mathcal A\) a prediction or action space. In many regression problems \(\mathcal A=\mathcal Y=\mathbb R\). In classification, \(\mathcal Y=\{1,\ldots,K\}\) while \(\mathcal A\) may be the simplex
\[
\Delta_K=\left\{p\in\mathbb R^K:p_k\geq0,\ \sum_{k=1}^K p_k=1\right\}
\]
of predicted class probabilities.

A data-generating distribution is a probability measure
\[
P \quad \text{on} \quad \mathcal X\times\mathcal Y.
\]
A training sample is usually modeled as independent draws
\[
S_n=((X_1,Y_1),\ldots,(X_n,Y_n)),\qquad (X_i,Y_i)\sim P.
\]
The iid assumption is a mathematical assumption about how the sample is produced, not an automatic property of every dataset.

A hypothesis class is a set \(\mathcal F\) of candidate functions
\[
f:\mathcal X\to\mathcal A.
\]
A loss function is a map
\[
\ell:\mathcal A\times\mathcal Y\to[0,\infty)
\]
where \(\ell(f(x),y)\) measures the cost of predicting \(f(x)\) when the observed output is \(y\). The population risk of \(f\) is
\[
\boxed{
R(f)=\mathbb E_{(X,Y)\sim P}\bigl[\ell(f(X),Y)\bigr].
}
\]
A learning algorithm is a map
\[
\mathcal L_n:(\mathcal X\times\mathcal Y)^n\to \mathcal F
\]
that takes the sample \(S_n\) and returns a learned function \(\widehat f=\mathcal L_n(S_n)\).

\paragraph{Core mechanism: the Bayes predictor for squared loss.}
Assume \(\mathcal Y=\mathcal A=\mathbb R\), and use squared loss
\[
\ell(a,y)=(a-y)^2.
\]
Among all measurable functions \(f:\mathcal X\to\mathbb R\), the risk-minimizing predictor is
\[
\boxed{
f^*(x)=\mathbb E[Y\mid X=x].
}
\]

\emph{Derivation.} Condition on \(X=x\). For any number \(a\),
\[
\mathbb E[(Y-a)^2\mid X=x]
=\mathbb E[(Y-\mu(x)+\mu(x)-a)^2\mid X=x],
\]
where \(\mu(x)=\mathbb E[Y\mid X=x]\). Expanding,
\[
\mathbb E[(Y-\mu(x))^2\mid X=x]
+2(\mu(x)-a)\mathbb E[Y-\mu(x)\mid X=x]
+(\mu(x)-a)^2.
\]
The middle term is zero by the definition of \(\mu(x)\). Thus
\[
\mathbb E[(Y-a)^2\mid X=x]
=\operatorname{Var}(Y\mid X=x)+(\mu(x)-a)^2,
\]
which is minimized at \(a=\mu(x)\). Since the total risk averages this conditional risk over \(X\), the pointwise minimizer gives the global minimizer.

\paragraph{Concrete example.}
Let \(X\in\{0,1\}\). Suppose
\[
P(Y=0\mid X=0)=P(Y=2\mid X=0)=1/2
\]
and
\[
P(Y=1\mid X=1)=1/4,\qquad P(Y=5\mid X=1)=3/4.
\]
For squared loss,
\[
f^*(0)=\mathbb E[Y\mid X=0]=1,
\]
and
\[
f^*(1)=\mathbb E[Y\mid X=1]=1(1/4)+5(3/4)=4.
\]
The best squared-loss predictor is not necessarily an observed response value for every input; it is the conditional mean.

\paragraph{Computational neuroscience example.}
Let \(X\) be a stimulus orientation and \(Y\) be a spike count in a \(100\) ms window. A tuning-curve model is a function
\[
f:\mathcal X\to\mathbb R
\]
that predicts mean response from stimulus. Under squared loss, the ideal function is \(f^*(x)=\mathbb E[Y\mid X=x]\), the conditional mean firing rate. Under a Poisson likelihood, the function may instead predict a nonnegative rate \(\lambda(x)\) whose likelihood is judged by spike-count probabilities. The biological statement is modest: the function summarizes trial-averaged response statistics under an experimental condition. It is not the full mechanism of spiking.

\paragraph{AI and machine-learning example.}
In image classification, \(X\) may be an image tensor and \(Y\in\{1,\ldots,K\}\) a label. A neural network often outputs
\[
f_\theta(x)\in\Delta_K,
\]
a vector of class probabilities. With cross-entropy loss
\[
\ell(f_\theta(x),y)=-\log (f_\theta(x))_y,
\]
the risk is expected negative log probability assigned to the true class. This is the same formal pattern as the tuning-curve example: choose a function class, choose a loss, and minimize expected loss.

\paragraph{Common misconceptions and failure modes.}
First, a training set is not the same object as the data distribution \(P\). The sample is finite; the distribution is the population model that generates future examples. Second, the target of learning depends on the loss. Squared loss targets conditional means, absolute loss targets conditional medians, and classification losses target class probabilities or decision rules. Third, a hypothesis class is an assumption. Restricting to linear functions, smooth functions, convolutional networks, or sparse models changes what can be learned.

\paragraph{Exercises.}
\begin{enumerate}[label=\textbf{31.1.\arabic*.}, leftmargin=*]
\item \textbf{Mechanical.} For a problem with \(x\in\mathbb R^{20}\), \(y\in\{0,1\}\), predicted probability \(a\in[0,1]\), and binary cross-entropy loss, identify \(\mathcal X\), \(\mathcal Y\), \(\mathcal A\), \(f\), and \(\ell\).
\item \textbf{Proof or derivation.} Repeat the squared-loss derivation and show that the conditional risk decomposes as \(\operatorname{Var}(Y\mid X=x)+(\mathbb E[Y\mid X=x]-a)^2\).
\item \textbf{Numerical/computational.} Let \(P(Y=0\mid X=a)=0.8\), \(P(Y=10\mid X=a)=0.2\). Compute the squared-loss Bayes prediction at \(X=a\).
\item \textbf{Interpretation.} In a neural decoding problem, \(X\) is a population spike-count vector and \(Y\) is a reach direction. Explain what distribution \(P(X,Y)\) represents.
\item \textbf{Misconception/debugging.} A student says, "The best predictor is just the most likely output." Explain why this can be false for squared loss and give a one-input example.
\end{enumerate}

\subsection{31.2 Empirical risk minimization}

\paragraph{Guiding question.}
If population risk depends on the unknown distribution \(P\), what objective can we actually minimize from data?

\paragraph{Beginner-facing intuition.}
The population risk \(R(f)\) is the quantity we care about, but it averages over examples we have not seen. Empirical risk minimization replaces the population distribution by the empirical distribution that puts mass \(1/n\) on each training example. The resulting objective is computable.

This replacement is powerful and dangerous. It makes learning possible, but it also makes overfitting possible: a function can do well on the sample because it adapted to noise or accidental patterns in that sample.

\paragraph{Formal definitions.}
Given a training sample \(S_n=((x_i,y_i))_{i=1}^n\), the empirical risk of \(f\in\mathcal F\) is
\[
\boxed{
\widehat R_n(f)=\frac{1}{n}\sum_{i=1}^n \ell(f(x_i),y_i).
}
\]
An empirical risk minimizer is any function
\[
\boxed{
\widehat f\in\operatorname*{arg\,min}_{f\in\mathcal F}\widehat R_n(f).
}
\]
If the minimum is not attained, algorithms often return an approximate minimizer. If \(\mathcal F=\{f_\theta:\theta\in\Theta\}\), empirical risk minimization becomes parameter optimization:
\[
\widehat \theta\in\operatorname*{arg\,min}_{\theta\in\Theta}
\frac{1}{n}\sum_{i=1}^n \ell(f_\theta(x_i),y_i).
\]

\paragraph{Likelihood-loss link.}
Many losses are negative log-likelihoods. Suppose a model assigns conditional density or mass
\[
p_\theta(y\mid x)
\]
to output \(y\) given input \(x\). The conditional log-likelihood of the sample is
\[
\sum_{i=1}^n \log p_\theta(y_i\mid x_i).
\]
Maximizing this is equivalent to minimizing
\[
\frac{1}{n}\sum_{i=1}^n -\log p_\theta(y_i\mid x_i).
\]
Thus negative log-likelihood is a loss. Gaussian noise gives squared loss; Bernoulli and categorical likelihoods give cross-entropy losses; Poisson likelihoods give count-model losses.

\paragraph{Core derivation: least squares as ERM and maximum likelihood.}
Let \(x_i\in\mathbb R^d\), \(y_i\in\mathbb R\), and use the linear model
\[
f_w(x)=w^\top x,\qquad w\in\mathbb R^d.
\]
With squared loss, empirical risk is
\[
\widehat R_n(w)=\frac{1}{n}\sum_{i=1}^n (w^\top x_i-y_i)^2.
\]
Put the rows \(x_i^\top\) into a design matrix \(X\in\mathbb R^{n\times d}\) and outputs into \(y\in\mathbb R^n\). Then
\[
\widehat R_n(w)=\frac{1}{n}\|Xw-y\|_2^2.
\]
Differentiating with respect to \(w\),
\[
\nabla_w \widehat R_n(w)=\frac{2}{n}X^\top(Xw-y).
\]
Setting the gradient to zero gives the normal equations
\[
\boxed{X^\top X\,\widehat w=X^\top y.}
\]
If \(X^\top X\) is invertible,
\[
\widehat w=(X^\top X)^{-1}X^\top y.
\]

The same objective follows from a Gaussian noise model
\[
Y_i\mid X_i=x_i\sim\mathcal N(w^\top x_i,\sigma^2),
\]
independently across \(i\). Up to constants independent of \(w\), the negative log-likelihood is
\[
\frac{1}{2\sigma^2}\sum_{i=1}^n (y_i-w^\top x_i)^2,
\]
so maximizing Gaussian likelihood and minimizing squared empirical risk choose the same \(w\).

\paragraph{Concrete example.}
Consider three data points for a one-parameter model \(f_a(x)=a x\):
\[
(x_i,y_i)=(1,1),(2,2),(3,2).
\]
For \(a=1\),
\[
\widehat R_3(f_1)=\frac{1}{3}\bigl[(1-1)^2+(2-2)^2+(3-2)^2\bigr]=\frac{1}{3}.
\]
For \(a=0.8\),
\[
\widehat R_3(f_{0.8})
=\frac{1}{3}\bigl[(0.8-1)^2+(1.6-2)^2+(2.4-2)^2\bigr]
=\frac{0.36}{3}=0.12.
\]
The empirical risk prefers \(a=0.8\) among these two candidates, even though neither candidate is guaranteed to minimize population risk.

\paragraph{Computational neuroscience example.}
An encoding model may predict firing rate from stimulus features:
\[
\lambda_\theta(x)=\exp(\theta^\top x).
\]
For spike counts \(y_i\), a Poisson likelihood gives empirical negative log-likelihood
\[
\frac{1}{n}\sum_{i=1}^n
\left(\lambda_\theta(x_i)-y_i\log \lambda_\theta(x_i)\right)
\]
up to constants independent of \(\theta\). Fitting this model is ERM with a loss matched to count data. A common failure mode is to minimize squared error on highly variable counts when a count likelihood would better reflect the support and variance structure.

\paragraph{AI and machine-learning example.}
Deep networks are usually trained by ERM with stochastic gradient methods. A minibatch objective
\[
\frac{1}{B}\sum_{i\in\mathcal B}\ell(f_\theta(x_i),y_i)
\]
estimates the full empirical risk gradient. The optimizer does not see \(P\) directly; it sees sampled examples and their losses. Validation data are held out to estimate whether empirical improvement transfers beyond the training set.

\paragraph{Common misconceptions and failure modes.}
First, ERM is not automatically generalization. It only minimizes the average loss on the observed sample. Second, the empirical minimizer may not be unique; many functions can have zero training loss. Third, likelihood and loss are linked only after a probabilistic noise model is specified. Fourth, optimization error matters: an algorithm may fail to find the ERM solution even if it is well defined.

\paragraph{Exercises.}
\begin{enumerate}[label=\textbf{31.2.\arabic*.}, leftmargin=*]
\item \textbf{Mechanical.} Write the empirical risk for a dataset \((x_i,y_i)_{i=1}^n\) using absolute loss \(\ell(a,y)=|a-y|\).
\item \textbf{Proof or derivation.} Derive the normal equations for least squares by differentiating \(\|Xw-y\|_2^2\) using the matrix-calculus convention from Chapter 7.
\item \textbf{Numerical/computational.} For data \((1,2),(2,2),(3,4)\), compute the empirical squared loss for \(f(x)=x\) and \(g(x)=1.2x\).
\item \textbf{Interpretation.} Explain why the negative log-likelihood for a Poisson spike-count model is a loss function.
\item \textbf{Misconception/debugging.} A model has zero training error and high validation error. Give two possible explanations in terms of ERM, hypothesis class size, or data leakage.
\end{enumerate}

\subsection{31.3 Generalization}

\paragraph{Guiding question.}
Why can a function with low training loss perform poorly on new examples, and what mathematical assumptions let us control this gap?

\paragraph{Beginner-facing intuition.}
Generalization is the difference between fitting the examples we saw and predicting examples drawn from the same population. Training loss is visible. Test loss is the quantity we want to estimate. A learning method generalizes when low empirical risk is evidence of low population risk.

The core reason this is hard is that the learned function depends on the same data used to evaluate training loss. A very flexible function class can adapt to noise, making training loss overly optimistic.

\paragraph{Formal definitions.}
For a fixed function \(f\), the generalization gap is
\[
R(f)-\widehat R_n(f).
\]
For a learned function \(\widehat f=\mathcal L_n(S_n)\), both terms are random because \(\widehat f\) depends on the sample. A train/test split estimates this gap by fitting \(\widehat f\) on training data and computing an independent test average
\[
\widehat R_{\mathrm{test}}(\widehat f)
=\frac{1}{m}\sum_{j=1}^m \ell(\widehat f(x_j^{\mathrm{test}}),y_j^{\mathrm{test}}).
\]

For squared-error regression, a useful pointwise bias-variance identity is the following. Suppose a random training set produces a random predictor \(\widehat f\), and for a fixed input \(x\),
\[
Y=f^*(x)+\varepsilon,\qquad \mathbb E[\varepsilon\mid X=x]=0,\qquad \operatorname{Var}(\varepsilon\mid X=x)=\sigma^2(x).
\]
Then
\[
\mathbb E\bigl[(\widehat f(x)-Y)^2\bigr]
=
\underbrace{\sigma^2(x)}_{\text{irreducible noise}}
+
\underbrace{\left(\mathbb E[\widehat f(x)]-f^*(x)\right)^2}_{\text{squared bias}}
+
\underbrace{\mathbb E\left[(\widehat f(x)-\mathbb E[\widehat f(x)])^2\right]}_{\text{variance}}.
\]
High bias means the method systematically misses the target. High variance means the fitted function changes strongly across samples.

\paragraph{Core theorem: a finite-class generalization bound.}
Assume \(\mathcal F\) is finite and that for every \(f\in\mathcal F\),
\[
0\leq \ell(f(X),Y)\leq 1.
\]
If \(S_n\) is iid from \(P\), then with probability at least \(1-\delta\), every \(f\in\mathcal F\) satisfies
\[
\boxed{
\left|R(f)-\widehat R_n(f)\right|
\leq
\sqrt{\frac{\log(2|\mathcal F|/\delta)}{2n}}.
}
\]

\emph{Derivation.} Fix one \(f\). The random losses
\[
Z_i=\ell(f(X_i),Y_i)
\]
are iid in \([0,1]\), with \(\mathbb E[Z_i]=R(f)\) and sample average \(\widehat R_n(f)\). Hoeffding's inequality gives
\[
P\left(|\widehat R_n(f)-R(f)|>\epsilon\right)
\leq 2e^{-2n\epsilon^2}.
\]
To make this hold for every \(f\in\mathcal F\), use the union bound:
\[
P\left(\exists f\in\mathcal F:\ |\widehat R_n(f)-R(f)|>\epsilon\right)
\leq 2|\mathcal F|e^{-2n\epsilon^2}.
\]
Set the right-hand side equal to \(\delta\) and solve for \(\epsilon\):
\[
2|\mathcal F|e^{-2n\epsilon^2}=\delta
\quad\Longrightarrow\quad
\epsilon=\sqrt{\frac{\log(2|\mathcal F|/\delta)}{2n}}.
\]
Thus, with probability at least \(1-\delta\), no function in the finite class has empirical risk far from population risk by more than this amount.

\paragraph{Concrete example.}
Let \(n=1000\), \(|\mathcal F|=100\), and \(\delta=0.05\). Then
\[
\sqrt{\frac{\log(2|\mathcal F|/\delta)}{2n}}
=
\sqrt{\frac{\log(4000)}{2000}}
\approx
\sqrt{0.00415}
\approx 0.064.
\]
For bounded losses, the bound says that all \(100\) candidate functions have empirical risks within about \(0.064\) of their population risks with probability at least \(0.95\). The bound becomes smaller with more data and larger with a larger class.

\paragraph{Computational neuroscience example.}
Suppose a decoder maps population activity to movement direction. If training and test trials are iid from the same task distribution, held-out test error estimates future performance on that task. If test trials come from a new day, a new behavioral context, or a shifted electrode population, the iid assumption may fail. Then a small training error can overstate real decoding performance. In neural data analysis, trial correlations, repeated stimulus blocks, and preprocessing leakage can all make test estimates too optimistic.

\paragraph{AI and machine-learning example.}
Large neural networks can have enough parameters to fit random labels, yet they often generalize on structured data. Classical finite-class bounds explain the basic role of sample size and complexity, but they can be numerically loose for modern networks. Practical generalization analysis therefore combines validation, regularization, architecture bias, optimization behavior, data augmentation, and distribution-shift tests.

\paragraph{Common misconceptions and failure modes.}
First, test error is not a magical truth; it is another estimate and needs independent test data. Second, a bound that is loose is not useless: it can still teach how sample size, confidence, and class size enter. Third, high training accuracy is not evidence of learning unless compared with validation or test performance. Fourth, generalization claims are always relative to a target distribution; performance can drop under distribution shift.

\paragraph{Exercises.}
\begin{enumerate}[label=\textbf{31.3.\arabic*.}, leftmargin=*]
\item \textbf{Mechanical.} Define training risk, test risk, population risk, and generalization gap for a learned function \(\widehat f\).
\item \textbf{Proof or derivation.} Fill in the algebra that solves \(2|\mathcal F|e^{-2n\epsilon^2}=\delta\) for \(\epsilon\).
\item \textbf{Numerical/computational.} Compute the finite-class bound for \(n=500\), \(|\mathcal F|=20\), and \(\delta=0.1\).
\item \textbf{Interpretation.} In the bias-variance identity, what kind of modeling choice tends to increase bias, and what kind tends to increase variance?
\item \textbf{Misconception/debugging.} A student reports test accuracy after tuning hyperparameters repeatedly on the same test set. Explain why the reported number may no longer estimate true test performance.
\end{enumerate}

\subsection{31.4 No-free-lunch and inductive bias}

\paragraph{Guiding question.}
Why is some assumption about the world necessary for learning from finite data?

\paragraph{Beginner-facing intuition.}
Learning needs a reason that observed examples say something about unobserved examples. Smoothness, locality, sparsity, low dimension, stationarity, compositionality, symmetry, and causal structure are examples of such reasons. Without any assumption, every continuation of the training data is possible, and no method has a universal advantage.

An inductive bias is the preference a learning method uses to choose among functions that agree with the data or nearly agree with the data. Bias is not automatically bad. It is the price of generalization.

\paragraph{Formal definitions.}
An inductive bias can enter through:
\[
\begin{array}{c}
\text{hypothesis class } \mathcal F,\quad
\text{loss } \ell,\quad
\text{regularizer } \Omega,\\[0.25em]
\text{optimizer},\quad
\text{data representation}.
\end{array}
\]
For example, choosing linear functions assumes the target is well approximated by a hyperplane. Choosing a convolutional network assumes local spatial structure and weight sharing are useful. Choosing an RBF kernel assumes smoothness in the chosen metric.

A simple no-free-lunch setting is binary classification on a finite input space \(\mathcal X\). A target function is any map
\[
t:\mathcal X\to\{0,1\}.
\]
If all target functions are considered equally possible, then labels at unseen inputs are independent fair bits after conditioning on the training labels.

\paragraph{Core theorem: no advantage on an unconstrained unseen label.}
Let \(\mathcal X\) be finite, and suppose the target function \(T:\mathcal X\to\{0,1\}\) is drawn uniformly from all \(2^{|\mathcal X|}\) binary labelings. A learner observes labels on a proper subset \(S\subset\mathcal X\). For any unseen point \(x_0\notin S\), any prediction rule based on the observed labels has probability \(1/2\) of predicting \(T(x_0)\) correctly.

\emph{Proof.} Conditioning on the labels observed on \(S\) rules out all target functions that disagree with those labels. For the unseen point \(x_0\), exactly half of the remaining target functions have \(T(x_0)=0\), and exactly half have \(T(x_0)=1\). Therefore
\[
P(T(x_0)=1\mid \text{observed labels})=P(T(x_0)=0\mid \text{observed labels})=1/2.
\]
Any learner must choose either \(0\), \(1\), or a randomized mixture. Its success probability is \(1/2\). The training labels do not help because the prior over target functions contains no relation between seen and unseen labels.

\paragraph{Concrete example.}
Let \(\mathcal X=\{a,b,c\}\), and suppose training observes \(t(a)=0\), \(t(b)=1\), leaving \(c\) unseen. Among all target functions agreeing with the training data, there are exactly two:
\[
t_0(a)=0,\ t_0(b)=1,\ t_0(c)=0,
\]
and
\[
t_1(a)=0,\ t_1(b)=1,\ t_1(c)=1.
\]
Without an assumption linking \(c\) to \(a\) or \(b\), predicting \(0\) and predicting \(1\) are equally justified on average.

\paragraph{Computational neuroscience example.}
A smooth tuning-curve model assumes that nearby stimulus orientations have similar expected responses. This is an inductive bias. It allows interpolation from measured orientations to unmeasured orientations. If the neuron's response were an arbitrary independent number at every orientation, smooth interpolation would fail. The smoothness assumption is biologically plausible in many sensory systems, but it is still an assumption to be checked against data.

\paragraph{AI and machine-learning example.}
Convolutional networks encode locality and translation equivariance: a feature detected in one image region can be useful elsewhere. Graph neural networks encode relational message passing along edges. Transformers encode token-wise interactions with shared attention machinery. Kernel methods encode similarity through \(k(x,x')\). These biases let models generalize from finite data by treating some inputs, features, or transformations as related.

\paragraph{Common misconceptions and failure modes.}
First, no-free-lunch does not say learning is impossible. It says learning is impossible uniformly over all possible worlds without assumptions. Second, "unbiased" is not automatically good in the broad learning sense; a method with no useful preference cannot extrapolate. Third, a powerful model still has biases, often through architecture, optimization, data augmentation, and training protocol. Fourth, an inductive bias can be wrong: a smooth model can oversmooth sharp thresholds, and a translation-equivariant model can fail when absolute position matters.

\paragraph{Exercises.}
\begin{enumerate}[label=\textbf{31.4.\arabic*.}, leftmargin=*]
\item \textbf{Mechanical.} List three sources of inductive bias in a learning pipeline and give one example of each.
\item \textbf{Proof or derivation.} Generalize the no-free-lunch unseen-point argument from binary labels to \(K\) equally likely labels.
\item \textbf{Numerical/computational.} For \(|\mathcal X|=5\) binary inputs and labels observed on \(3\) inputs, how many target functions remain possible? How many assign label \(1\) to a particular unseen input?
\item \textbf{Interpretation.} Explain why smoothness is a useful bias for fitting a sensory tuning curve but may fail for a discontinuous decision rule.
\item \textbf{Misconception/debugging.} A student says, "A large neural network has no inductive bias because it can represent almost anything." Identify at least two remaining biases in architecture or training.
\end{enumerate}

\paragraph{Closing bridge.}
This chapter framed learning as statistical function approximation: define the population objective, compute a sample objective, and ask whether the fitted function generalizes. Chapter 32 now specializes this abstract setup to linear regression, logistic regression, generalized linear models, and margin-based classifiers. Chapter 33 then shows how regularization and kernels express inductive bias more explicitly, and how Gaussian processes turn functions themselves into random objects.

\clearpage
